% ============================================================
%  Propuesta de Proyecto: Agente de Trading MCTS — TSLA
% ============================================================
\documentclass[11pt, a4paper]{article}

% ── Encoding & fonts ─────────────────────────────────────────
\usepackage[T1]{fontenc}
\usepackage[utf8]{inputenc}
\usepackage[spanish]{babel}
\usepackage{lmodern}
\usepackage{microtype}

% ── Page layout ──────────────────────────────────────────────
\usepackage[top=2.2cm, bottom=2.2cm, left=2.4cm, right=2.4cm]{geometry}

% ── Color palette ────────────────────────────────────────────
\usepackage[dvipsnames, svgnames, x11names]{xcolor}
\definecolor{bg}{HTML}{FFFFFF}
\definecolor{card}{HTML}{F6F8FA}
\definecolor{acc1}{HTML}{0969DA}   % azul
\definecolor{acc2}{HTML}{1A7F37}   % verde
\definecolor{acc3}{HTML}{CF222E}   % naranja/rojo
\definecolor{acc4}{HTML}{8250DF}   % violeta
\definecolor{txt}{HTML}{1F2328}
\definecolor{mut}{HTML}{57606A}

% ── Graphics & diagrams ──────────────────────────────────────
\usepackage{tikz}
\usetikzlibrary{arrows.meta, shapes.geometric, shapes.misc,
                positioning, fit, calc, backgrounds,
                decorations.pathreplacing, matrix}
\usepackage{pgfplots}
\pgfplotsset{compat=1.17}

% ── Boxes & layout helpers ───────────────────────────────────
\usepackage[most]{tcolorbox}
\tcbuselibrary{breakable, skins, theorems}
\usepackage{booktabs}
\usepackage{tabularx}
\usepackage{array}
\usepackage{multicol}
\usepackage{enumitem}
\usepackage{amsmath, amssymb}
\usepackage{mathtools}
\usepackage{fontawesome5}

% ── Hyperlinks ───────────────────────────────────────────────
\usepackage[colorlinks=true, linkcolor=acc1, urlcolor=acc1]{hyperref}

% ── Page colors ──────────────────────────────────────────────
\pagecolor{bg}
\color{txt}

% ── tcolorbox styles ─────────────────────────────────────────
\tcbset{
  darkcard/.style={
    enhanced, breakable,
    colback=card, colframe=#1, coltitle=white,
    boxrule=0.8pt, arc=4pt,
    left=8pt, right=8pt, top=6pt, bottom=6pt,
    fonttitle=\bfseries\color{#1},
  },
  metricbox/.style={
    enhanced,
    colback=#1!10!card, colframe=#1,
    boxrule=1pt, arc=5pt,
    halign=center, valign=center,
  }
}

% ── Section style ─────────────────────────────────────────────
\usepackage{titlesec}
\titleformat{\section}
  {\large\bfseries\color{acc1}}
  {\thesection}{0.7em}{}[{\color{acc1}\titlerule[0.5pt]}]

\titleformat{\subsection}
  {\normalsize\bfseries\color{acc2}}
  {\thesubsection}{0.6em}{}

% ── Helpers ──────────────────────────────────────────────────
\newcommand{\highlight}[2][acc1]{\textcolor{#1}{\textbf{#2}}}
\setlength{\parskip}{4pt}
\setlength{\parindent}{0pt}

% =============================================================
\begin{document}

% ─── PORTADA ─────────────────────────────────────────────────
\begin{titlepage}
  \centering
  \vspace*{1.5cm}

  % Accent line top
  {\color{acc1}\rule{\linewidth}{1.5pt}}
  \vspace{0.4cm}

  {\footnotesize\color{mut}\MakeUppercase{Propuesta de Proyecto — Marzo 2026}}
  \vspace{0.8cm}

  {\Huge\bfseries\color{txt}
    Agente de Trading MCTS\\[0.3em]
    Estocástico para TSLA}

  \vspace{0.5cm}
  {\color{acc1}\rule{\linewidth}{1.5pt}}
  \vspace{1.2cm}

  % Badge pills
  \begin{tikzpicture}
    \foreach \label/\col [count=\i] in {
      MDP/acc1, MCTS/acc2, GBM/acc3, Backtesting/acc4}{
      \node[rounded corners=6pt,
            fill=\col!20!card,
            draw=\col, line width=0.8pt,
            text=\col, font=\bfseries\small,
            minimum width=2.4cm, minimum height=0.65cm]
        at ({(\i-1)*3.0},0) {\label};
    }
  \end{tikzpicture}

  \vspace{1.4cm}

  % Mini-abstract box
  \begin{tcolorbox}[darkcard=acc1, width=0.88\linewidth,
    title={Resumen ejecutivo}]
    Se propone diseñar un \highlight{agente de trading} capaz de tomar
    decisiones óptimas de compra/venta sobre la acción de Tesla (TSLA)
    modelando el mercado como un \highlight{Proceso de Decisión de Markov}
    y empleando \highlight{Monte Carlo Tree Search} (el algoritmo de
    AlphaGo) para buscar la acción de mayor valor esperado.
    El entorno estocástico se simula con \highlight{Movimiento Browniano
    Geométrico} y la estrategia se valida contra \emph{Buy \& Hold}
    mediante backtesting out-of-sample sobre datos de Yahoo Finance.
  \end{tcolorbox}

  \vfill
  {\small\color{mut}
    \faDatabase\; Datos: Yahoo Finance (TSLA)\quad
    \faCode\; Algoritmo: MCTS + GBM\quad
    \faChartLine\; Benchmark: Buy \& Hold}
\end{titlepage}

% ─── 1. FORMULACIÓN MDP ──────────────────────────────────────
\section{Formulación: Proceso de Decisión de Markov}

El problema de trading se modela como el MDP $(S, A, P, R)$.
Cada instante $t$ el agente observa el estado del portafolio,
elige una acción y recibe una recompensa en forma de retorno logarítmico.

% ── Diagrama ciclo MDP ──────────────────────────────────────
\begin{center}
\begin{tikzpicture}[
  node distance=2.6cm,
  every node/.style={font=\small},
  box/.style={rounded corners=6pt, draw=#1, fill=#1!15!card,
              text=txt, minimum width=2.8cm, minimum height=1.0cm,
              align=center, font=\bfseries\small, text=txt},
  arr/.style={-{Stealth[length=7pt]}, thick, color=#1},
]
  \node[box=acc1] (S)  {Estado $s_t$\\
    {\footnotesize\color{mut}$(P_t, C_t, H_t)$}};
  \node[box=acc2, right=of S] (A)  {Acción $a_t$\\
    {\footnotesize\color{mut}$\{-k\%,\,0,\,+k\%\}$}};
  \node[box=acc3, below=1.8cm of A] (E)  {Entorno\\
    {\footnotesize\color{mut}Mercado / GBM}};
  \node[box=acc4, below=1.8cm of S] (R)  {Recompensa $R_t$\\
    {\footnotesize\color{mut}$\ln(V_t/V_{t-1})$}};

  \draw[arr=acc2] (S) -- node[above, text=mut, font=\scriptsize]{decide} (A);
  \draw[arr=acc3] (A) -- node[right, text=mut, font=\scriptsize]{aplica} (E);
  \draw[arr=acc4] (E) -- node[below, text=mut, font=\scriptsize]{genera} (R);
  \draw[arr=acc1] (R) -- node[left,  text=mut, font=\scriptsize]{actualiza} (S);
\end{tikzpicture}
\end{center}

\begin{tcolorbox}[darkcard=acc1,
  title={Componentes del MDP}]
\renewcommand{\arraystretch}{1.4}
\begin{tabularx}{\linewidth}{@{} l l X @{}}
  \toprule
  \color{acc1}\textbf{Símbolo} &
  \color{acc1}\textbf{Componente} &
  \color{acc1}\textbf{Definición} \\
  \midrule
  $s_t = (P_t, C_t, H_t)$ & Estado &
    Precio de cierre, capital líquido y acciones en cartera.
    El valor total es $V_t = C_t + H_t \cdot P_t$. \\
  $a_t \in \{-k\%,\,0,\,+k\%\}$ & Acción &
    Vender, mantener o comprar el $k$\,\% del capital disponible. \\
  $R_t = \ln\!\left(\dfrac{V_t}{V_{t-1}}\right)$ & Recompensa &
    Retorno logarítmico del portafolio entre dos instantes consecutivos. \\
  \bottomrule
\end{tabularx}
\end{tcolorbox}

% ─── 2. MOTOR MCTS ───────────────────────────────────────────
\section{Motor MCTS — Las cuatro fases}

El agente itera las siguientes fases para construir un árbol de búsqueda
asimétrico y seleccionar la acción óptima en cada instante:

% ── Diagrama árbol MCTS ─────────────────────────────────────
\begin{center}
\begin{tikzpicture}[
  level 1/.style={sibling distance=52mm},
  level 2/.style={sibling distance=22mm},
  every node/.style={circle, draw, minimum size=0.9cm,
                     inner sep=1pt, font=\scriptsize\bfseries},
  edge from parent/.style={draw, -stealth, color=mut, thick},
  sel/.style   ={draw=acc1, fill=acc1!20!card, text=txt},
  exp/.style   ={draw=acc2, fill=acc2!20!card, text=txt},
  roll/.style  ={draw=acc3, fill=acc3!20!card, text=txt, dashed},
  back/.style  ={draw=acc4, fill=acc4!20!card, text=txt},
  norm/.style  ={draw=mut,  fill=card,          text=mut},
]
% root
\node[sel] (root) {$s_0$}
  child { node[norm] (l1) {$s_{1}$} }
  child { node[sel, edge from parent/.style={draw=acc1,thick,-stealth}]
            (l2) {$s_{2}^*$}
    child { node[norm] (l21) {$s_{3}$} }
    child { node[exp,  edge from parent/.style={draw=acc2,thick,-stealth}]
              (l22) {$s_{4}$}
      child[missing] {}
    }
    child { node[norm] (l23) {$s_{5}$} }
  }
  child { node[norm] (l3) {$s_{6}$} };

% rollout chain
\node[roll, right=1.2cm of l22] (r1) {$\tau_1$};
\node[roll, right=1.0cm of r1]  (r2) {$\tau_2$};
\node[roll, right=1.0cm of r2]  (r3) {$\tau_T$};
\draw[-stealth, dashed, color=acc3, thick] (l22) -- (r1);
\draw[-stealth, dashed, color=acc3, thick] (r1) -- (r2);
\draw[dotted, color=acc3, thick] (r2) -- (r3);

% backprop arrow
\draw[-stealth, color=acc4, very thick, bend left=40]
  (r3.north) to node[draw=none, fill=none, rectangle, midway,
    text=acc4, font=\scriptsize\bfseries] {$R$}
  (root.east);

% Phase labels
\node[draw=none, fill=none, rectangle, text=acc1,
      font=\scriptsize\bfseries, above=0.2cm of root]
  {\faSearch\; Selección (UCT)};
\node[draw=none, fill=none, rectangle, text=acc2,
      font=\scriptsize\bfseries, below left=0.1cm and -0.3cm of l22]
  {\faPlus\; Expansión};
\node[draw=none, fill=none, rectangle, text=acc3,
      font=\scriptsize\bfseries, above=0.2cm of r2]
  {\faDice\; Rollout (GBM)};
\node[draw=none, fill=none, rectangle, text=acc4,
      font=\scriptsize\bfseries, right=0.3cm of root]
  {\faArrowUp\; Backprop};
\end{tikzpicture}
\end{center}

\vspace{0.3em}
\begin{tcolorbox}[darkcard=acc2, title={Las cuatro fases en detalle}]

\textbf{\color{acc1}\faSearch\; Selección.}
Desde la raíz $s_0$, el algoritmo desciende por el árbol eligiendo en
cada nodo la acción que maximiza la política UCT:
\[
  a^* = \arg\max_{a \in A}
    \left[
      \underbrace{Q(s,a)}_{\text{explotación}}
      +\; c\sqrt{\dfrac{\ln N(s)}{n(s,a)}}
    \right]
\]
donde $Q(s,a)$ es la recompensa media observada, $N(s)$ el total de
visitas al nodo padre, $n(s,a)$ las visitas a esa acción y $c$ el
coeficiente de exploración.

\medskip
\textbf{\color{acc2}\faPlus\; Expansión.}
Si el nodo seleccionado tiene acciones no exploradas, se instancia un
nuevo nodo hijo $s_{t+1}$ aplicando una de ellas.

\medskip
\textbf{\color{acc3}\faDice\; Simulación (Rollout).}
Desde el nuevo nodo se simula el mercado hasta el horizonte $T$ usando
un Movimiento Browniano Geométrico (ver sección siguiente).

\medskip
\textbf{\color{acc4}\faArrowUp\; Retropropagación.}
La recompensa acumulada $R$ se propaga hacia la raíz actualizando
contadores y valores:
\[
  n(s,a) \leftarrow n(s,a)+1,\qquad
  Q(s,a) \leftarrow Q(s,a) + \frac{R - Q(s,a)}{n(s,a)}
\]
\end{tcolorbox}

% ─── 3. MODELO ESTOCÁSTICO ───────────────────────────────────
\section{Modelo estocástico del entorno — GBM}

Durante el rollout el precio futuro de TSLA se simula como un
\highlight[acc3]{Movimiento Browniano Geométrico}:
\[
  P_{t+1} = P_t \cdot \exp\!\left[
    \left(\mu - \tfrac{\sigma^2}{2}\right)\Delta t
    + \sigma\sqrt{\Delta t}\; Z
  \right],\qquad Z \sim \mathcal{N}(0,1)
\]

\begin{center}
\begin{tikzpicture}
\begin{axis}[
  width=14cm, height=6cm,
  axis background/.style={fill=card},
  tick style={color=mut},
  ticklabel style={color=mut, font=\scriptsize},
  xlabel={\color{mut}\footnotesize Días de trading},
  ylabel={\color{mut}\footnotesize Precio TSLA (\$)},
  title={\color{txt}\small Trayectorias GBM desde $P_0=250\,\$$
         ($\mu=0.03\%$\,día, $\sigma=1.8\%$\,día)},
  every axis title/.style={at={(0.5,1.02)}},
  ymajorgrids=true, grid style={color=card!200!bg, line width=0.3pt},
  xmin=0, xmax=252, ymin=80, ymax=600,
  axis line style={color=mut!60},
  clip=false,
]
% Background paths (faint)
\addplot[acc1!50!card, opacity=0.5, line width=0.7pt, samples=253, domain=0:252, no marks]
  {250*exp(0.0003*x + 0.018*sqrt(x)*0.8)};
\addplot[acc2!50!card, opacity=0.5, line width=0.7pt, samples=253, domain=0:252, no marks]
  {250*exp(0.0003*x - 0.018*sqrt(x)*0.3)};
\addplot[acc4!50!card, opacity=0.5, line width=0.7pt, samples=253, domain=0:252, no marks]
  {250*exp(0.0003*x + 0.018*sqrt(x)*1.3)};
\addplot[acc3!50!card, opacity=0.5, line width=0.7pt, samples=253, domain=0:252, no marks]
  {250*exp(0.0003*x - 0.018*sqrt(x)*1.0)};
\addplot[acc1!30!card, opacity=0.4, line width=0.7pt, samples=253, domain=0:252, no marks]
  {250*exp(0.0003*x + 0.018*sqrt(x)*0.2)};
\addplot[acc2!30!card, opacity=0.4, line width=0.7pt, samples=253, domain=0:252, no marks]
  {250*exp(0.0003*x - 0.018*sqrt(x)*1.6)};
\addplot[acc4!30!card, opacity=0.4, line width=0.7pt, samples=253, domain=0:252, no marks]
  {250*exp(0.0003*x + 0.018*sqrt(x)*1.8)};
\addplot[mut!60!card, opacity=0.35, line width=0.7pt, samples=253, domain=0:252, no marks]
  {250*exp(0.0003*x - 0.018*sqrt(x)*2.0)};
% Two highlighted paths
\addplot[acc2, line width=2pt, samples=253, domain=0:252, no marks]
  {250*exp(0.0006*x + 0.018*sqrt(x)*0.4)};
\addplot[acc3, line width=2pt, samples=253, domain=0:252, no marks]
  {250*exp(-0.0001*x - 0.018*sqrt(x)*0.55)};
\addplot[color=mut, dashed, line width=0.8pt] coordinates {(0,250)(252,250)};
\node[text=mut, font=\scriptsize] at (axis cs:240, 262) {$P_0$};
\end{axis}
\end{tikzpicture}
\end{center}

\begin{tcolorbox}[darkcard=acc3, title={Parámetros del modelo}]
\renewcommand{\arraystretch}{1.35}
\begin{tabularx}{\linewidth}{@{} c l X l @{}}
  \toprule
  \color{acc3}\textbf{Parámetro} &
  \color{acc3}\textbf{Significado} &
  \color{acc3}\textbf{Descripción} &
  \color{acc3}\textbf{Valor típico} \\
  \midrule
  $\mu$        & Deriva        & Retorno medio diario histórico de TSLA
                               & $\approx 0.03\%$ \\
  $\sigma$     & Volatilidad   & Desviación estándar diaria de los retornos
                               & $\approx 1.8\%$ \\
  $\Delta t$   & Paso temporal & Granularidad de la simulación
                               & 1 día \\
  $T$          & Horizonte     & Duración máxima del rollout
                               & 252 días (1 año) \\
  \bottomrule
\end{tabularx}
\end{tcolorbox}

% ─── 4. EVALUACIÓN ───────────────────────────────────────────
\section{Evaluación empírica — Backtesting}

El modelo se validará con datos históricos de TSLA obtenidos de
\highlight{Yahoo Finance} usando validación cruzada temporal
(\emph{walk-forward}). Se compararán las siguientes métricas frente
a la estrategia pasiva \emph{Buy \& Hold}:

\begin{center}
\begin{tikzpicture}[node distance=0cm]
  \node[rounded corners=6pt, draw=acc1, fill=acc1!12!card,
        text width=6cm, align=center, inner sep=12pt] (roi) {
    {\color{acc1}\Large\bfseries ROI}\\[4pt]
    {\color{mut}\footnotesize Rendimiento acumulado}\\[6pt]
    $\displaystyle\frac{V_T - V_0}{V_0}$\\[6pt]
    {\color{mut}\scriptsize Mide el retorno total sobre el periodo}
  };
  \node[rounded corners=6pt, draw=acc2, fill=acc2!12!card,
        text width=6cm, align=center, inner sep=12pt,
        right=1cm of roi] (sharpe) {
    {\color{acc2}\Large\bfseries Sharpe}\\[4pt]
    {\color{mut}\footnotesize Retorno ajustado por riesgo}\\[6pt]
    $\displaystyle S = \frac{\mathbb{E}[R_p - R_f]}{\sigma_p}$\\[6pt]
    {\color{mut}\scriptsize Penaliza estrategias de alta varianza}
  };
\end{tikzpicture}
\end{center}

\vspace{0.5em}

\begin{tcolorbox}[darkcard=acc4, title={Esquema de validación}]
\begin{enumerate}[leftmargin=*, itemsep=3pt]
  \item Descargar serie histórica de TSLA (mínimo 3 años).
  \item Dividir en ventanas de entrenamiento y test (\emph{walk-forward}).
  \item Ejecutar el agente MCTS en cada ventana de test.
  \item Calcular ROI y Sharpe para MCTS y \emph{Buy \& Hold}.
  \item Comparar distribuciones de retornos con test estadístico.
\end{enumerate}
\end{tcolorbox}

% ─── 5. EXTENSIONES ──────────────────────────────────────────
\section{Extensiones posibles}

\begin{tcolorbox}[darkcard=mut]
\begin{itemize}[leftmargin=*, itemsep=4pt]
  \item \highlight[acc1]{Estado enriquecido}: añadir indicadores técnicos
    (RSI, MACD, Bandas de Bollinger) al vector de estado $s_t$.
  \item \highlight[acc2]{Función de valor aprendida}: sustituir el rollout
    GBM por una red neuronal (estilo AlphaZero) entrenada sobre datos
    históricos.
  \item \highlight[acc3]{Múltiples activos}: extender el espacio de acciones
    a una cartera de activos correlacionados.
  \item \highlight[acc4]{Costes de transacción}: incorporar comisiones y
    \emph{slippage} en la función de recompensa.
\end{itemize}
\end{tcolorbox}

\vfill
{\footnotesize\color{mut}\centering
  Propuesta elaborada para evaluación del grupo · Marzo 2026 ·
  Datos: Yahoo Finance · Algoritmo: MCTS + GBM}

\end{document}
